\chapter*{卷二\quad 岁计}
\phantomsection\addcontentsline{toc}{chapter}{卷二\quad 岁计}
\markboth{卷二\quad 岁计}{卷二\quad 岁计}

卷二十片，乃全书中数字密度最高之一卷，所记为常例定额而非特定年份之收支。此卷释文最称完整，校记却最为棘手：诸片数字虽俱全，其所指涉之具体口径往往欠明，每能读出确数，而无从确知此数所从出之统计边界。

\begin{jian}{简十}
岁计口，凡三十二万一千。匠六万，工十一万，余为居者。

凡三口出一工。工者不入匠籍。匠者不入工籍。
\end{jian}
\chuchu{原件残片 37 号，长 20.5 cm，金属，字口较浅。}
\xiaojì{「三十二万一千」较简十四所称三十一万二千为多。此片字口较浅，与他片不类，疑为另一次清点，或属草稿；同卷之中两数并存而不相核校，为他片所无。\par{}「匠六万，工十一万」，两项相加为十七万，余十五万一千为居者，三项之和恰好为三十二万一千，与本片总数相合。「凡三口出一工」为征工之法；若按三口出一工推算，十一万工当出于三十三万人，与全数不合，可见「工」与「口」并非同一套编制口径，此点下编第六章另论。}

\begin{jian}{简十一}
匠籍别为一册。工籍别为一册。居者不入册，以口计。

册各有吏，不相知。
\end{jian}
\chuchu{原件残片 39 号，长 14.8 cm，陶瓷。}
\xiaojì{「居者不入册」，说明城中人口占比最大之一部分不在任何名册之上。此点对下编第六章的推算影响甚巨：所有残存的分类数字皆系于十七万有籍者，而十五万一千居者仅余一总数。「册各有吏，不相知」乃一项关键之制度设计：三册分掌、互不相通，则任何一名吏员皆无从掌握全城实口。此制与简八「板与册不同者以板为正」合观，足见城中对记录权限之分割一以贯之。}

\begin{jian}{简十二}
粟支九年，水支六年，气支三年。

三物之支，同时而计，不同时而尽。

尝有一岁，粟不足三月。减自上始，居者日一升，工如故。
\end{jian}
\chuchu{原件残片 41 号，长 9.4 cm，金属，全书之中尺寸偏小的一片。}
\xiaojì{三数并列而无主语，所支者为库存总量，抑或每人每年之定额，无从骤断。若为库存总量，则可与简十四所记口数、简十五所记日给相权，进而推出仓容；若作定额解，则全无着落。本书取下编第六章所论之前一说。「同时而计，不同时而尽」为全书至要之一语：三物之中气最先绝，水次之，粟最后；若三者确为存量，则此城在理论上存在一明确之「断气之日」，且该时点远在撤离之前，与简二十七所记地热之衰相参，撤离时点之推算详下编第七章。\par{}末段记一次粮荒：某年存粟仅支三月，处置乃「减自上始」，居者每日减至一升，工匠照旧。一升为定额之半；按简十五「人日二升」计，居者口粮减半而工匠不减，正与简十五「不足则减，减自上始」之序相合，即先减居者，次减工，最后减匠。此条将「减」字落于实事，证明其确属一次实际执行之削减而非虚文，受削减者涵盖全城十五万一千居者中之大多数；至于事发何年，本片未记。}

\begin{jian}{简十三}
新入者，先居外舍三月，乃入环。入环之日，给籍，给衣，给器。

器不继，以旧器补。旧器有记，记其旧主。
\end{jian}
\chuchu{原件残片 43 号，长 17.1 cm，金属。}
\xiaojì{「新入者」，指自地上归来之人。此类人在城中占比甚小，然制度仍为其专设程序，可见往来未曾完全断绝；三月之期亦与他处所记隔离期相合。「旧器有记，记其旧主」，谓凡补给之旧器上皆保留原主标记。此一细节与全书其余处所见制度气质相背：一个以「名可削」为常制之地，却在器物上留存旧主之记，此点下编第八章有专论。}

\begin{figure}[htbp]
\centering
\begin{tikzpicture}[x=1mm,y=1mm]
  \useasboundingbox (-8,-16) rectangle (76,30);
  % 片身
  \fill[shade] (0,0) rectangle (56,18);
  \draw[ink,line width=0.5pt] (0,0) rectangle (56,18);
  % 行界与字界
  \foreach \x in {7,14,21,28,35,42,49}{\draw[rulegrey,line width=0.15pt] (\x,0) -- (\x,18);}
  \draw[rulegrey,line width=0.3pt] (0,9) -- (56,9);
  % 第一行：计口三十一万二千
  \foreach \i/\c in {1/计,2/口,3/三,4/十,5/一,6/万,7/二,8/千}{
    \node[font=\scriptsize] at (\i*7-3.5,13.5) {\c};}
  % 第二行：实居十一万□千
  \foreach \i/\c in {1/实,2/居,3/十,4/一,5/万,6/九}{
    \node[font=\scriptsize] at (\i*7-3.5,4.5) {\c};}
  \node[font=\scriptsize] at (49-3.5,4.5) {千};
  % 蚀去（不在其上压线，只填浅灰并以细边标出）
  \fill[faintgrey] (42,0) rectangle (49,9);
  \draw[rulegrey,line width=0.15pt] (42,0) rectangle (49,9);
  \node[font=\scriptsize,text=pencil] at (49-3.5,4.5) {□};
  % 尺寸
  \draw[<->,rulegrey,line width=0.3pt] (0,-3) -- (56,-3);
  \node[font=\scriptsize,text=pencil,anchor=north] at (28,-4) {18.4 cm（原件实测）};
  \draw[<->,rulegrey,line width=0.3pt] (60,0) -- (60,18);
  \node[font=\scriptsize,text=pencil,anchor=west,align=left] at (61,9) {1.7 cm\\两行};
  \node[font=\scriptsize,text=pencil,anchor=south west] at (0,18.6) {残片 26 号};
  \node[font=\scriptsize,text=pencil,anchor=south east] at (56,18.6) {（今字对录）};
\end{tikzpicture}
\caption{简十四（残片 26 号）形制与刻写行款摹本}\label{fig:slip}
\figsource{据原件照片与一九六三年编目卡重绘。原件所刻为小篆，此处照释文体例改作今字，非原字形；灰色一格为蚀去之处，即校记所补「百」字的位置。}
\end{figure}

\begin{jian}{简十四}
计口三十一万二千，实居十一万九千。

粟支九年，水支六年，气支三年。
\end{jian}
\chuchu{原件残片 26 号，长 18.4 cm，金属，正面有朱色残痕。编号与卷一次序不合，系一九六三年二次编目所致。}
\xiaojì{「十一万九千」下残片蚀去一处。胶片本作「十一万九千六百」，残片无可核对，故其字不取；本书所用之数仍取十一万九千，与简二十、简二十四相同。\par{}此片为全书最关紧要之一片：额定与实居并列，差额十九万三千，即空置六成二。第二行三数与简十二全同，可知简十二乃此片之下半，二者或为同一长条文书分刻于两片者；若此说成立，则简十二「同时而计，不同时而尽」一句，当紧接此片之后。}

\begin{jian}{简十五}
月给粟，人日二升。匠倍之。外工又倍之。

不足则减，减自上始。
\end{jian}
\chuchu{原件残片 45 号，长 13.6 cm，金属。}
\xiaojì{按人日二升计，十一万九千人日需二千三百八十石，岁需八十六万八千石；与简十二「粟支九年」相权，则仓容当在七百万石以上，此数构成本书下编第六章推算之基础。「不足则减，减自上始」为分配原则，谓先减居者，后减工，最后减匠；此句虽简短，却为全书唯一涉及配给优先次序之明文。}

\begin{jian}{简十六}
工籍：捣者三千，负者一千二百，修者八百，余为杂役。

捣者不入环，月给别。
\end{jian}
\chuchu{原件残片 47 号，长 16.9 cm，陶瓷。}
\xiaojì{「捣者三千」，三千之数为全书仅见，与他处工匠之数俱不相合；「捣」字在此处无宾语，所捣为何物，本片未言。「捣者不入环，月给别」，谓此三千人不进入居住环，口粮另行发放。此制与简五「外工」相类而微异：外工明确居于垣外，捣者则只言不入环而未明其居处，二者是否同一群体，残简无据可断。}

\begin{jian}{简十七}
定额：人日三合，多者赏，少者罚。罚不过三，过三则移籍。

赏以粟，罚以役。

十四年，捣者十一人连日不足，罚役三日，不移籍。其一人再犯，移籍。
\end{jian}
\chuchu{原件残片 49 号，长 12.3 cm，金属。}
\xiaojì{「三合」为容积单位，所量当指成品；此片与简十六文义相接，可断为同属一册。「赏以粟，罚以役」谓奖惩一律以劳役与口粮结算，不涉他物，可见城中激励机制颇为单一。「过三则移籍」之移籍即改隶他册；参以简十一「册各有吏，不相知」，则移籍实为跨系统之调动，非一名吏员所能独断，唯移往何籍本片未言。\par{}末段记一次具体处分：十四年，捣者十一人连日未足定额，各罚役三日而不改籍，其中一人再犯始予移籍。此为全书唯一可考见「罚不过三」如何施行之实录：十一人同时不足，足见当日症结多在定额本身，然处分仍照律施行，仅再犯者被移籍；至于是何一人再犯，本片亦未记。}

\begin{jian}{简十八}
凡籍，岁一易。易籍之日，旧籍焚。焚者二人，见者三人，皆署名。

焚余之灰，埋于垣外。

十九年，一册不焚。焚者自陈，见者三人同辞。罚焚者，册终不焚。
\end{jian}
\chuchu{原件残片 51 号，长 15.8 cm，金属。}
\xiaojì{旧籍岁焚，乃全书最大之一处史料空洞：一九六三年所得者，皆为焚余或未及焚毁之件而非全帙，下编所论皆须以此为前提。「焚者二人，见者三人，皆署名」为焚籍之成规，每次焚籍皆须五人留名，而现存三十七片中绝无此五人之署名，可见所焚之籍与所署之签俱不在此。「焚余之灰，埋于垣外」与简六「限外百步不得置物」相抵，二禁相权，当以葬灰为特例。\par{}末段记十九年一桩特例：一册未焚，司焚者自陈其事，在场三人供辞一致，结果罚司焚者而该册终未焚毁。「终不焚」三字未附解释，既可解作未及补焚，亦可解作有意存留，本片不偏执一端。此条实关涉本书之成立：倘若所有旧籍皆依章程焚尽，一九六三年所能觅得者便唯有灰烬；今人之所以仍得见此三十七片，正因该制度在严密运转之中，终究存在一两桩「不焚」之例外。}

\begin{jian}{简十九}
事故之数：兔乱二万二千。

逃亡者，籍除，名削。
\end{jian}
\chuchu{原件残片 53 号，长 19.2 cm，金属，右端有灼痕。}
\xiaojì{「兔乱」二字，「兔」字全书仅两见，此为其一；乱事死者之数，他片俱未及，乃本书独有之材料。二万二千为城中事变死者之总数，与简十四「计口」及「实居」之差额并非同一笔账，详下编第六章。「逃亡者，籍除，名削」为逃亡处置之法：逃者既除籍削名，便不在任何册籍之内，此点与简二十五「无名者不得登」相接，构成全书第二处重大矛盾。}
